\chapter{结论}\label{chap:conclusion}

\section{主要工作总结}

本研究围绕全细胞分子动力学模拟中的长程静电求解问题，构建了一个面向膜体系代理模型的研究原型，并提出了基于随机批量采样与高斯和分解的长程静电方法（Random Batch Sum-of-Gaussians, RBSOG）。围绕这一方法，本文完成了算法实现、实验流程搭建、结果统计与论文资产自动生成等一整套工作。

本研究的主要工作包括：

（1）\textbf{构建了面向膜体系的分子动力学模拟原型系统。}该系统实现了直接求和参考解、PPPM基线方法和RBSOG方法，并提供了命令行入口、模拟配置接口、结果记录、统计分析与图表导出流程。

（2）\textbf{建立了系统化的基准测试体系。}该体系从计算性能、统计方差和结构稳定性三个维度对PPPM和RBSOG进行比较，评测指标覆盖平均每步时间、压力方差、温度统计、能量漂移、面积/脂质波动与膜厚代理量漂移等关键量。

（3）\textbf{通过真实实验明确了RBSOG当前阶段的边界。}在本文重新串行执行并核验的真实实验中，RBSOG在当前纯Python原型中的绝对运行速度显著慢于PPPM：三种子主基准下约慢6.71倍，十种子扩展统计下约慢\RBSOGSlowdownVsPPPM 倍；与此同时，RBSOG的压力方差均值也高于PPPM。

（4）\textbf{构建了可复现的自动化实验与论文生成流程。}该流程能够支持单次模拟、多随机种子基准对比、批量大小扫描、JIT消融实验和结构稳定性分析，并自动生成图表、表格和LaTeX资产。

\section{创新点总结}

本研究的主要创新点包括：

（1）\textbf{提出了RBSOG长程静电求解框架。}将高斯和分解与随机批量采样结合到长程静电求解中，为面向大规模分子动力学模拟的低通信路线提供了一个可实现原型。

（2）\textbf{建立了兼顾计算效率、统计方差和结构稳定性的综合评测体系。}除传统的性能指标外，还纳入了压力方差、结构代理量、效应量分析与Pareto推荐方法。

（3）\textbf{构建了可复现的实验与论文联动流程。}实验目录、自动生成图表与论文正文之间形成了明确的数据链路，降低了结果复核和后续迭代的成本。

\section{贡献与意义}

本文的贡献主要体现在以下几个方面。

\subsection{理论与方法学意义}

本文证明了“高斯和分解 + 随机批量采样”这一组合路线可以被迁移到长程静电求解场景中，并在分子动力学研究原型内稳定运行。这为后续继续探索低通信型长程静电算法提供了一个具体起点\cite{jin2020random,guo2025sum}。

\subsection{工程与研究流程意义}

本文同时建立了一套可复现实验流程，使单次运行、多种子比较、批量扫描、JIT消融和论文资产生成能够被统一管理。这一部分工作虽然不直接等同于算法创新，但对于后续高质量迭代具有明确价值\cite{feig2023current}。

\section{最终结论}

基于本文的真实实验结果，可以得出一个清晰而克制的结论：\textbf{RBSOG在当前纯Python原型中尚未在性能和统计稳定性上超过PPPM。}这一定性在三种子主基准、十种子扩展统计、批量扫描与JIT消融实验中保持一致。

不过，这一结论并不意味着RBSOG路线无效。更准确的理解是：当前工作已经完成了“方法可实现、流程可复现、问题可量化”的原型验证阶段，但尚未完成“方法已优于成熟基线”的证据积累。JIT消融实验表明，仅实现层优化就可带来约\JITSpeedup 倍的提速，这说明当前原型距离路线的性能上限仍有显著距离。

因此，本文最重要的产出不是证明RBSOG已经胜出，而是以真实实验结果明确了其当前边界、暴露了后续工作重点，并为下一阶段的工程优化、规模扫描与更真实体系验证提供了可靠起点。
